\section{La radiografía de tórax y los datos}

\subsection{Cómo se forma una radiografía}

Una radiografía es una \textbf{proyección}: registra en un plano la atenuación
acumulada de un haz de rayos X que atraviesa todo el tórax~\cite{prince2015}. El
haz se produce en un tubo en el que electrones acelerados por una diferencia de
potencial (kVp) frenan contra un ánodo; la corriente y el tiempo de exposición
(mAs) fijan la cantidad de fotones.

Al atravesar un material de coeficiente de atenuación $\mu$ y espesor $x$, la
intensidad decae según la ley de Beer-Lambert:
\begin{equation}
I = I_0 \, e^{-\mu x}.
\end{equation}
El coeficiente $\mu$ crece con la densidad y el número atómico. El hueso y el
mediastino atenúan mucho; el aire del parénquima pulmonar, casi nada. La imagen se
representa de modo que las regiones de \textbf{alta atenuación se ven claras}
(hueso, corazón, mediastino) y las de \textbf{baja atenuación, oscuras} (campos
pulmonares). Parte de los fotones se dispersa y degrada el contraste; por eso se
usa una rejilla antidispersión.

El detector convierte la intensidad que llega en una señal digital cuantizada. Los
dos conjuntos usan tecnologías distintas:

\begin{itemize}
    \item \textbf{Radiografía computarizada (CR)}, en Montgomery: una placa de
    fósforo fotoestimulable almacena la imagen latente, que luego se lee con un
    láser y se digitaliza.
    \item \textbf{Radiografía digital directa (DR)}, en Shenzhen: un panel plano
    convierte los rayos X en señal digital sin paso intermedio.
\end{itemize}

La respuesta de contraste y la estructura del ruido difieren entre ambas, de modo
que un mismo procesamiento no produce el mismo resultado en los dos orígenes. Esa
diferencia es la condición de comparación del proyecto.

\subsection{Protocolo posteroanterior (PA)}

Todas las imágenes son radiografías de tórax en proyección PA. La tabla resume el
protocolo estándar; los conjuntos no documentan los parámetros de cada exposición.

\begin{table}[H]
\centering
\small
\caption{Elementos del protocolo estándar de la radiografía de tórax PA}
\begin{tabular}{@{}p{4cm}p{9.6cm}@{}}
\toprule
\textbf{Elemento} & \textbf{Descripción y motivo} \\
\midrule
Dirección del haz & De la espalda (posterior) al pecho (anterior), con el tórax
apoyado en el detector \\
Posición & De pie, con los hombros rotados hacia adelante para retirar las
escápulas de los campos pulmonares \\
Distancia foco-detector & Alrededor de 180 cm: reduce la magnificación y la
distorsión geométrica \\
Respiración & Inspiración profunda sostenida: expande los pulmones y desciende el
diafragma \\
Técnica & Alto kilovoltaje, para atravesar el mediastino y equilibrar el contraste
entre hueso y pulmón \\
Rejilla & Antidispersión, para conservar el contraste \\
Ventaja frente a AP & El corazón queda cerca del detector y se magnifica menos \\
\bottomrule
\end{tabular}
\end{table}

En la imagen se distinguen las clavículas, las costillas (bandas casi horizontales
que cruzan los campos), la columna, el corazón y el mediastino en el centro, los
hemidiafragmas abajo y los dos \textbf{campos pulmonares}, que son la región que
segmenta este proyecto.

\subsection{Conjuntos de datos}

Se usaron dos conjuntos públicos de la U.S. National Library of
Medicine~\cite{jaeger2014}: \textbf{Montgomery County Chest X-ray Set} (Estados
Unidos, \FinNMC{} imágenes) y \textbf{Shenzhen Hospital Chest X-ray Set} (China,
\FinNSH{} imágenes), en total \FinNImagenes{} radiografías, de las cuales
\FinNConMascara{} tienen máscara de referencia. La etiqueta (con o sin hallazgos de
tuberculosis) proviene de las lecturas clínicas que acompañan cada imagen.

\input{final_datos}

\paragraph{Por qué dos conjuntos.} En la formulación se compararon cuatro
candidatos (JSRT, Montgomery, Shenzhen y CheXmask). Ningún conjunto individual
permite responder la pregunta, porque todos provienen de un solo sitio y la
pregunta pide mediciones reproducibles \emph{entre} sitios. JSRT se descartó por
riesgo de acceso, y CheXmask por razones metodológicas: sus máscaras las generó una
red neuronal, y usarlas como referencia invalidaría la comparación contra un
trazado humano.

\paragraph{Máscaras de referencia.} Montgomery entrega dos archivos booleanos por
imagen, uno por pulmón, trazados bajo supervisión de radiólogo~\cite{candemir2014}.
Las de Shenzhen provienen del Instituto Politécnico de Kiev~\cite{stirenko2018}:
un archivo de 512 por 512 píxeles con bordes suavizados, que se reescala por vecino
más cercano y se binariza en 128. Las dos procedencias siguen convenciones
distintas: \textbf{las de Montgomery excluyen el corazón y los grandes vasos; las
de Shenzhen los delimitan}. Es un sesgo sistemático entre orígenes que no puede
corregirse, y es la principal limitación del estudio.

\begin{figure}[H]
\centering
\includegraphics[width=\textwidth]{galeria_entrenamiento.png}
\caption{Galería de la cohorte de entrenamiento: las dos clases y los dos
orígenes, los extremos de resolución, un paciente pediátrico y un caso marcado
para revisión. Selección aleatoria con semilla fija dentro de cada criterio.}
\end{figure}

\subsection{Caracterización de la digitalización (OE1)}
\label{sec:caracterizacion}

La caracterización se hizo sobre las \FinNImagenes{} imágenes, sin descartar
ninguna, y corrigió varios supuestos de la documentación y de la propia
formulación:

\begin{table}[H]
\centering
\small
\caption{Lo que se suponía en la formulación y lo que se encontró}
\begin{tabular}{@{}p{3.2cm}p{5cm}p{5.6cm}@{}}
\toprule
\textbf{Propiedad} & \textbf{Supuesto inicial} & \textbf{Hallazgo} \\
\midrule
Profundidad & Montgomery de 12 bits & \FinBitsMC{} bits efectivos en Montgomery y
\FinBitsSH{} en Shenzhen (archivos \texttt{uint8}) \\
Escala invertida & Shenzhen invertida & \FinInvertidas{} imágenes invertidas: no se
aplicó corrección \\
Canales & Shenzhen en RGB & Tres canales idénticos, verificados uno a uno; se toma
un canal \\
Tamaño & Montgomery homogéneo & Montgomery: siempre 4020 por 4892, con relleno negro
de hasta el 67\,\% de la imagen. Shenzhen: 525 tamaños distintos en 662 imágenes \\
Orientación & No considerada & Ambos mezclan orientaciones: 97 verticales y 41
horizontales en Montgomery; 368 y 294 en Shenzhen \\
Archivos anómalos & Nueve archivos pequeños & Son los nueve pacientes menores de
dos años: el campo se ajusta al paciente \\
\bottomrule
\end{tabular}
\end{table}

\begin{figure}[H]
\centering
\includegraphics[width=\textwidth]{dimensiones.png}
\caption{Dimensiones por origen en la cohorte de entrenamiento. Montgomery entrega
siempre el mismo tamaño de archivo; Shenzhen lo ajusta al paciente.}
\end{figure}

\paragraph{Variabilidad de intensidad.} La mediana de intensidad es de 111,0 en
Montgomery y 173,5 en Shenzhen (en escala de 0 a 255), mientras que dentro de cada
origen las clases apenas difieren. \textbf{El efecto del sitio supera al de la
clase}: un descriptor de intensidad que separara las clases estaría separando
hospitales. Por eso toda métrica del proyecto se reporta estratificada por origen.

\begin{figure}[H]
\centering
\includegraphics[width=\textwidth]{intensidad.png}
\caption{Histogramas por imagen y resumen por origen en la cohorte de
entrenamiento. Los dos orígenes forman nubes casi disjuntas.}
\end{figure}

\subsection{Calidad y partición}

Una imagen es \emph{aceptada} si pasa los controles de lectura, dimensiones y
canales. Las imágenes sin máscara de referencia participan en la caracterización,
pero no en la evaluación contra la referencia. Ningún conjunto incluye
identificador de paciente: las lecturas clínicas permitieron demostrar dos pacientes
con varias imágenes, que se agrupan para que nunca crucen particiones.

La partición se fijó una sola vez (70/15/15 por paciente, estratificada por origen
y clase, semilla 20262) y se conservó en todos los laboratorios. La edad no se
estratificó, para no dejar celdas de uno o dos elementos; su reparto se verificó
después y es comparable entre particiones.

\input{final_particion}